% ECONOMICS_PROBLEM: {"id": "L74-446", "title": "Construction-productivity stagnation", "jel_code": "L74", "source_id": "OP-0462"}
% !TeX program = xelatex
\documentclass[11pt,a4paper]{article}
\usepackage[margin=23mm]{geometry}
\usepackage{fontspec}
\setmainfont{Latin Modern Roman}
\usepackage{amsmath,amssymb,mathtools}
\usepackage{xcolor,enumitem,booktabs,longtable,array,xurl,hyperref}
\definecolor{muted}{HTML}{53636C}
\hypersetup{colorlinks=true,urlcolor=blue,linkcolor=blue}
\setlength{\parindent}{0pt}
\setlength{\parskip}{5pt}
\newcommand{\R}{\mathbb R}
\newcommand{\E}{\mathbb E}
\newcommand{\N}{\mathbb N}
\newcommand{\ind}{\mathbf 1}
\newcommand{\Var}{\operatorname{Var}}
\newcommand{\Cov}{\operatorname{Cov}}
\newcommand{\supp}{\operatorname{supp}}
\DeclareMathOperator*{\argmin}{arg\,min}
\DeclareMathOperator*{\argmax}{arg\,max}
\newcommand{\entrymeta}[3]{{\sffamily\small\color{muted}\textbf{\texttt{#1}}\quad|\quad #2\par #3\par}}
\newcommand{\Problemref}[1]{\texttt{#1}}
\newcommand{\JELref}[2]{\texttt{#2} (JEL \texttt{#1})}
\begin{document}
\section*{Construction-productivity stagnation}
\textbf{Problem ID:} \texttt{L74-446}\quad \textbf{JEL:} \texttt{L74}\par
% BEGIN ECONOMICS STATEMENT
\label{problem:OP-0462}
{\sffamily\small\color{muted}Original subject group: Urban, housing and spatial economics\par}
\label{definitions:problem:30007744}
\textbf{Audit classification:} Published research agenda. Checked source identity and appropriate agenda/background support; expanded intervention, units, population, definedness and causal restrictions. Exact targets and variants are editorial.
\subsection*{Definitions and precise research target}
\textit{Editorial formalization.} Consider residential construction establishments in one country from 2000 through 2025. Let \(Q_{it}\) be quality-adjusted completed floor area, \(L_{it}\) labor hours, \(K_{it}\) capital services, and \(M_{it}\) materials for establishment \(i\) in year \(t\). For operating establishments with positive output and input aggregator, define productivity \(A_{it}=Q_{it}/F(L_{it},K_{it},M_{it})\), where the production aggregator \(F\) and quality adjustments are stated and estimated independently of the policy contrast. Let \(r\in\{0,1\}\) denote current land-use rules versus a specified streamlined regime. Let \(O_i(r)\) indicate operation at follow-up under regime \(r\). The survivor causal target is
\[\tau_A=E[\log A_{i,t+5}(1)-\log A_{i,t+5}(0)\mid O_i(1)=O_i(0)=1].\]
Assume the always-operating stratum has positive probability and integrable log contrasts. Its membership is latent: use justified principal-stratum assumptions or bounds, and report survival effects separately. Identify whether regulation changes productivity through project scale, establishment entry, standardization, or technology adoption, rather than merely changing output prices or building composition. Track subcontracting and unrecorded hours; aggregate industry productivity additionally requires weights and an entry treatment. Use credible rule changes with justified counterfactual trends and test sensitivity to alternative output deflators. Decompose aggregate stagnation only when within-firm changes, reallocation, and quality are separately identified. The housing survey explicitly requests research on construction-productivity causes and remedies. The source does not provide this exact establishment-level estimand; its target and measurement restrictions are an adapted specification.

Inputs are positive measured service flows and \(F(L,K,M)>0\) a predeclared production index; \(A\) is a physical-output productivity residual, distinct from revenue productivity. \(O_i(r)\in\{0,1\}\) indicates operation at year five, and the always-operating principal stratum has probability greater than zero. \(\tau_A\) is defined on latent potential survival membership and generally is not point identified even by randomized reform. Monotone survival \(O_i(1)\geq O_i(0)\), bounded \(\log A\), and assignment independence can give bounds; any tighter principal ignorability must be stated separately. Aggregate productivity additionally needs fixed quality units, establishment weights and explicit entry/exit accounting. Regulation may alter output mix and subcontracting, so a changing deflator cannot be silently interpreted as technological stagnation. All expectations use a stated baseline population and probability law. Identification means every admissible data-generating model with the same observed-data law yields the same displayed target; otherwise report the identified set. Exact equations, horizons and assumptions are editorial, rather than source-given canonical conjectures.
\subsection*{Variants and assumption changes}
\begin{enumerate}
\item \textbf{Editorial survival-bounds variant.} Drop principal ignorability; impose randomized reform, bounded log productivity and monotone survival. Bound survivor effects and report survival separately.
\item \textbf{Editorial industry variant.} Use total quality-adjusted output over an aggregate input index including entry/exit. Decompose within-establishment and reallocation changes with declared weights.
\end{enumerate}
\subsection*{References and source audit}
\textbf{1.} Explicit agenda on dynamic supply, construction productivity, and richer migration models; exact targets are editorial. Locator: Handbook of Regional and Urban Economics, Volume 6, Chapter 6; Section 7, pp.448–449. Full text or primary publication page inspected. URL: \url{https://realestate.wharton.upenn.edu/wp-content/uploads/2025/12/BaumSnow_Duranton_bc_2025-003.pdf}.
\begin{enumerate}\raggedright
\item\hypertarget{definitions:source:30007744:1}{} Nathaniel Baum-Snow; Gilles Duranton. 2025. \emph{Housing supply and housing affordability}. Handbook of Regional and Urban Economics, Volume 6, Chapter 6; Section 7, pp.448–449.
\url{https://realestate.wharton.upenn.edu/wp-content/uploads/2025/12/BaumSnow_Duranton_bc_2025-003.pdf}.
DOI: \url{https://doi.org/10.1016/bs.hesreg.2025.05.003}.
\end{enumerate}

\subsection*{Common conventions: Source mathematical conventions}

These conventions apply to each entry unless explicitly replaced.

\textbf{Models and identification.} A primitive model \(m\) specifies all probability laws, feasible choices, information, equilibrium and measurement rules used by its target. The admissible class \(\mathcal M\) is fixed independently of outcomes. If \(P_O(m)\) is its observable probability law and \(\tau(m)\) the target, its identified set at observable law \(P\) is
\[
 \mathcal I_\tau(P)=\{\tau(m):m\in\mathcal M,\ P_O(m)=P\}.
\]
Point identification means that this set is a singleton. An empty set signals incompatibility of data and assumptions. Coordinate bounds need not describe the joint set. Bounds are sharp only if every retained target is attainable or arbitrarily approachable by compatible models. Finite bounds require explicit restrictions on unobserved quantities; unrestricted confounding, hidden wealth, extrapolation or arbitrary equilibrium selection can yield uninformative sets.

\textbf{Causal interventions.} A potential outcome \(Y_i(p)\) is the outcome under a complete policy regime \(p\), including timing, financing, information and market spillovers. The target population and units are fixed before treatment. With interference, use the complete assignment vector or a justified exposure mapping; individual no-interference notation alone cannot identify an economy-wide intervention. A difference of mean potential outcomes does not require the cross-world joint distribution. Conditioning on post-treatment migration, survival, enrollment or employment changes the estimand and needs separate justification.

\textbf{Statistical statements.} An estimator, confidence set, forecast or test specifies a sampling law, model class and loss/coverage criterion. Uniform \((1-\alpha)\)-coverage means \(\inf_{m\in\mathcal M}\Pr_m\{\tau(m)\in C_n\}\geq1-\alpha\), or an explicitly asymptotic version. The whole parameter space trivially has coverage, so informativeness requires a stated width, risk or power objective. A forecasting improvement is relative to a named benchmark, information set and evaluation population.

\textbf{Equilibrium and optimization.} An equilibrium comprises feasible plans, optimality, beliefs where needed, budget constraints and all market/resource clearing. The primitives or a stated benchmark define each correspondence \(\mathcal E(p;m)\). Nonemptiness is assumed or proved, never inferred just from compact policy sets. A selected-equilibrium target uses a declared selection rule; a robust target quantifies over all permitted equilibria. Use suprema or infima unless attainment is established. An \(\varepsilon\)-optimal policy is feasible and within \(\varepsilon\) of the finite optimum value. Report an empty feasible set separately.

\textbf{Welfare and accounting.} Normative utility, interpersonal normalization, social weights, numeraire, discounting and the use of public receipts are primitives. Behavioral utility need not coincide with the welfare criterion. Household welfare includes ownership income and rents once; adding profits again to compensating variation generally double-counts them. A monetary equivalent is defined by an explicit transfer and price convention, with monotonicity and range conditions. Average effects, mechanism contributions, transfers and welfare losses are different targets. Infinite sums require convergence and dynamic feasible plans require terminal or no-Ponzi conditions.

\textbf{Source versions.} A working-paper date, revision date and journal publication year are distinguished. A DOI for a journal version is not automatically the DOI of an earlier manuscript. Locators refer to the stated version and printed pages where specified. A metadata-only check does not verify a claim about a full-text conclusion. Every entry's source subsection narrows the provenance claim accordingly.
% END ECONOMICS STATEMENT
\end{document}
